
We asked whether better SSD approximation at long context is accompanied by worse retrieval, and whether these two phenomena are causally related. The mechanism gate result establishes that they are not: MOHAWK's SSD mixer approximates LLaMA-3.1-8B attention matrices with decreasing normalized error as sequence length grows ($\beta = -0.368$, 1,920 measurements, both gate criteria passed with large margin at N $\leq$ 2,048). The approximation is not failing at long context. Any retrieval degradation in fully-trained MOHAWK-SSM models must therefore be attributed to the state update h\_t = $A\cdot h_{t-1} + B\cdot x_t$, which exponentially discounts early token positions regardless of how well the SSD weights are initialized --- a structural property of the bounded-state architecture, not a fixable training artifact.

Four contributions are reported: (1) mechanism gate confirmed --- normalized SSD error decreases with N ($\beta = -0.368$, 1,920 measurements); (2) architectural attribution --- retrieval degradation in MOHAWK-converted LLaMA-3.1-8B is a bounded-state property, not a distillation quality artifact; (3) validated evaluation infrastructure --- the first controlled MOHAWK-SSM versus LAWCAT comparison pipeline on fixed LLaMA-3.1-8B with LongBench v2 is implemented, tested (22/22 unit tests), and running; (4) ten engineering failure modes identified and resolved. The behavioral confirmation (H-E1, H-M2) awaits experiment completion.

